% ECONOMICS_PROBLEM: {"id": "C61-246", "title": "Finite-horizon drift control with discretionary stopping", "jel_code": "C61", "source_id": "OP-0568"}
% !TeX program = xelatex
\documentclass[11pt,a4paper]{article}
\usepackage[margin=23mm]{geometry}
\usepackage{fontspec}
\setmainfont{Latin Modern Roman}
\usepackage{amsmath,amssymb,mathtools}
\usepackage{xcolor,enumitem,booktabs,longtable,array,xurl,hyperref}
\definecolor{muted}{HTML}{53636C}
\hypersetup{colorlinks=true,urlcolor=blue,linkcolor=blue}
\setlength{\parindent}{0pt}
\setlength{\parskip}{5pt}
\newcommand{\R}{\mathbb R}
\newcommand{\E}{\mathbb E}
\newcommand{\N}{\mathbb N}
\newcommand{\ind}{\mathbf 1}
\newcommand{\Var}{\operatorname{Var}}
\newcommand{\Cov}{\operatorname{Cov}}
\newcommand{\supp}{\operatorname{supp}}
\DeclareMathOperator*{\argmin}{arg\,min}
\DeclareMathOperator*{\argmax}{arg\,max}
\newcommand{\entrymeta}[3]{{\sffamily\small\color{muted}\textbf{\texttt{#1}}\quad|\quad #2\par #3\par}}
\newcommand{\Problemref}[1]{\texttt{#1}}
\newcommand{\JELref}[2]{\texttt{#2} (JEL \texttt{#1})}
\begin{document}
\section*{Finite-horizon drift control with discretionary stopping}
\textbf{Problem ID:} \texttt{C61-246}\quad \textbf{JEL:} \texttt{C61}\par
% BEGIN ECONOMICS STATEMENT
\label{problem:OP-0568}
\label{beyond:SC-07}
\entrymeta{SC-07}{Stochastic control and dynamic optimization}{Research agenda formalization}
\subsection*{Definitions and assumptions}
Use the preceding control and stopping conventions. Fix a deterministic
maturity $H>0$. For a state $X_t=x$ at $t\in[0,H]$, let
\[
 V_H(t,x)=\inf_{u,\tau:\,t\leq\tau\leq H}
 \E\left[k(X_\tau)+\int_t^\tau(\psi(u_s)+c)\,ds\right].
\]
Immediate stopping gives $0\leq V_H(t,x)\leq k(x)$ and
$V_H(H,x)=k(x)$. Define the stopping set
$\mathcal S_H=\{(t,x):V_H(t,x)=k(x)\}$ and its complement as the
continuation set. The infimum is over the weak admissible systems just
defined, with the state initialized at $(t,x)$.
\subsection*{Formal question}
Characterize $V_H$, an optimal control and an optimal stopping time for
this coefficient/cost class. Determine conditions under which
\[
 \mathcal S_H=\{(t,x):|x|\leq b_H(t)\}
\]
for a moving boundary $b_H:[0,H]\to[0,\infty]$ with
$b_H(H)=+\infty$, including its regularity for $t<H$ and behavior as
$t\uparrow H$. Establish a verification theorem for the associated
obstacle equation
\[
 \min\left\{k-V_H,\ \partial_tV_H+
 \tfrac12\sigma^2\partial_{xx}V_H+
 \inf_{a\in\R}\{a\mu\partial_xV_H+\psi(a)\}+c\right\}=0.
\]
Where spatial differentiability holds, the candidate feedback is
$(\psi')^{-1}(-\mu(x)\partial_xV_H(t,x))$; its admissibility and
optimality must be proved.
\subsection*{Scope and status}
The interval form of the stopping set is a target to be justified under
specified assumptions, rather than a theorem asserted for all data.
The source proposes this finite-horizon extension and anticipates a moving
boundary; the formal obstacle problem above makes that agenda precise.
\subsection*{References for the source of the problem}
V\'aclav E. Bene\v{s}, Georgy Gaitsgori and Ioannis Karatzas,
\emph{Drift Control with Discretionary Stopping for a Diffusion},
arXiv:2401.10043v3, \S5, second research direction (finite maturity);
\S2, (A1)--(A7), equations (2.1)--(2.3), for the underlying model.
\url{https://arxiv.org/html/2401.10043v3#S5}.

\subsection*{Common conventions: Source mathematical conventions}

\textbf{Well-defined objects.}
All probability spaces and measurable variables are specified or inherited
from a local statistical model. Expectations used as finite targets require
the stated integrability. A decision rule or estimator is measurable with
respect to its available information; randomized rules may use an auxiliary
independent random variable. Norms without a subscript are Euclidean in
finite-dimensional spaces. An optimizer is requested only when attainment
is assured; otherwise the question concerns the optimal value and
$\varepsilon$-optimal admissible rules for every $\varepsilon>0$.

\textbf{Identification.}
A structural model $m\in\mathcal M$ implies an observable law
$P_m^{\rm obs}$ and target $\theta(m)$. The target is point identified on
$\mathcal M$ if
\[
 P_{m_1}^{\rm obs}=P_{m_2}^{\rm obs}
 \quad\Longrightarrow\quad\theta(m_1)=\theta(m_2)
 \qquad(m_1,m_2\in\mathcal M).
\]
When this fails, the sharp identified set is
\[
 \Theta(P;\mathcal M)=\{\theta(m):m\in\mathcal M,
                                  P_m^{\rm obs}=P\}.
\]
Specifying an intervention or estimand does not establish identification.
Positivity, exclusion, causal sufficiency, measurement restrictions, and
transport assumptions are substantive local inputs, never automatic
consequences of notation.

\textbf{Minimax risk and nontrivial inference.}
For a statistical experiment with data law $P^{(n)}$, target $\theta(P)$,
and nonnegative loss $\ell$, define
\[
 \mathcal R_n(\mathcal P,\ell)
   =\inf_{\widehat\theta_n}\sup_{P\in\mathcal P}
      \E_{P^{(n)}}\ell(\widehat\theta_n,\theta(P)).
\]
The infimum is over measurable estimators. A sharp rate is a sequence
$r_n$ for which $0<c\leq C<\infty$, independent of $n$, satisfies
$c r_n\leq\mathcal R_n\leq C r_n$ for all sufficiently large $n$.
Squared-error parametric risk is of order $n^{-1}$; the corresponding
root-mean-squared error is of order $n^{-1/2}$. These different conventions
are distinguished locally. A uniform asymptotic confidence region $C_n$
must satisfy
\[
 \liminf_{n\to\infty}\inf_{P\in\mathcal P}
     \Pr_{P^{(n)}}\{\theta(P)\in C_n\}\geq 1-\alpha,
 \qquad 0<\alpha<1,
\]
together with an informative diameter or expected-width requirement.
Coverage alone permits the vacuous whole parameter space. Testing questions
similarly impose both size control and power against a declared separated
alternative class.

\textbf{Binary causal baseline.}
When an entry explicitly invokes this baseline, one observes independent
copies of $O=(X,W,Y)$, with $W\in\{0,1\}$ and potential outcomes $Y(0),Y(1)$.
Assume consistency $Y=Y(W)$, no interference, conditional exchangeability
$(Y(0),Y(1))\perp W\mid X$, and overlap
$\epsilon\leq e(X)\leq1-\epsilon$ almost surely for a fixed
$0<\epsilon<1/2$, where
\[
 e(x)=\Pr(W=1\mid X=x),\qquad
 m_w(x)=\E[Y\mid W=w,X=x],\qquad
 \tau(x)=m_1(x)-m_0(x).
\]
These assumptions identify conditional mean effects almost everywhere
under the covariate law. Pointwise or uniform effects require appropriate
regularity and a specified support. Continuous treatment, longitudinal
data, adaptive experiments, and network interference require their own
assumptions in place of this baseline. Bounded outcomes, densities, and
smoothness are additional assumptions only where stated.

\textbf{Function classes.}
On $[0,1]^d$, $\mathcal H_d^s(L)$ is the isotropic Hölder ball of radius
$L>0$: put $k=\lceil s\rceil-1$ and $\gamma=s-k\in(0,1]$, bound derivatives
through order $k$, and require the order-$k$ derivatives to be
$\gamma$-Hölder with constant at most $L$. Thus integer orders use a
Lipschitz final derivative convention. The class
$\operatorname{BV}_d(L)$ consists of functions $f\in L^1((0,1)^d)$ whose
distributional derivative is a finite vector measure and for which
$\|f\|_{L^1}+|Df|((0,1)^d)\leq L$.
An $L^\infty$ bound or a particular pointwise version is imposed separately
when needed. Sparse and neural-network classes require a declared
dictionary or architecture and coefficient bounds.

An anisotropic H\"older class with declared block smoothness indices means
coordinatewise H\"older regularity in each block, uniformly in the other
blocks, with all corresponding derivative and H\"older bounds at most
the stated radius. Derivative orders and the integer-order convention in
each block are those just given. Mixed-derivative bounds are additional
restrictions only when explicitly stated.

\textbf{Decision and computational questions.}
Policy regret compares admissible feasible policies under the same law and
constraints. In computational entries, finite data and thresholds have
rational encodings; running time is measured in their total bit length.
Real-valued equilibrium search uses the explicitly declared algebraic or
FIXP encoding. An approximation ratio for maximization is written
$\operatorname{OPT}/\operatorname{ALG}$ and for minimization as
$\operatorname{ALG}/\operatorname{OPT}$, with zero optima and infeasible
instances handled locally. Historical approximation bounds are not
automatically current bounds.

\subsection*{Common conventions: Control and stopping conventions for SC-07--SC-08}

For the next two entries, retain the diffusion and cost assumptions of
Bene\v{s}, Gaitsgori and Karatzas, \S2: $\mu$ is even, positive and
differentiable on $(0,\infty)$; $\sigma$ is even, continuous and bounded
away from zero on $(\varepsilon,\infty)$ for every $\varepsilon>0$;
$\mu'(x)\sigma(x)^2=A\mu(x)^2$ for $x>0$, with $A\geq0$.
If $A=0$, $\sigma$ is bounded; if $A>0$, $\mu'$ is bounded away from
zero on each $(\varepsilon,\infty)$. The nonnegative, even functions
$k,\psi\in C^2(\R)$ are strictly convex, $\psi(0)=0$,
$\psi'(r)\to\infty$ as $r\to\infty$, and $c>0$.

An admissible control/state system is a weak solution, on a filtered
probability space carrying a Brownian motion $W$, with progressive real
control $u$ and continuous state $X$ satisfying
\[
 dX_s=u_s\mu(X_s)\,ds+\sigma(X_s)\,dW_s,
 \quad
 \int_0^t\bigl(|u_s\mu(X_s)|+\sigma(X_s)^2\bigr)ds<\infty
 \quad\text{a.s. for every }t<\infty.
\]
Stopping times are taken in the usual augmentation of the state
filtration. Infima range over all such weak systems and their associated
stopping times; nonnegative costs define extended expectations.
This is a single decision maker's optimization problem.
% END ECONOMICS STATEMENT
\end{document}
